\documentclass[11pt,a4paper]{article}
\usepackage[margin=27mm,headheight=15pt]{geometry}
\usepackage[T1]{fontenc}
\usepackage[utf8]{inputenc}
\usepackage{lmodern,microtype,amsmath,amssymb,amsthm,mathtools}
\usepackage{booktabs,longtable,array,enumitem,xcolor,fancyhdr}
\usepackage{listings}
\usepackage[colorlinks=true,linkcolor=blue!45!black,citecolor=blue!45!black,urlcolor=blue!50!black,breaklinks=true]{hyperref}
\usepackage{xurl}
\definecolor{ink}{RGB}{24,43,65}
\definecolor{accent}{RGB}{19,91,105}
\pagestyle{fancy}\fancyhf{}
\fancyhead[L]{\small\textcolor{ink}{Stable observables in iPRA}}
\fancyhead[R]{\small Research report}
\fancyfoot[C]{\thepage}
\setlength{\parskip}{4pt}
\setlength{\parindent}{0pt}
\setlist{topsep=4pt,itemsep=3pt,parsep=0pt}
\setlength{\emergencystretch}{2em}
\newtheorem{theorem}{Theorem}[section]
\newtheorem{lemma}[theorem]{Lemma}
\newtheorem{proposition}[theorem]{Proposition}
\newtheorem{corollary}[theorem]{Corollary}
\theoremstyle{definition}
\newtheorem{definition}[theorem]{Definition}
\newtheorem{example}[theorem]{Example}
\theoremstyle{remark}\newtheorem{remark}[theorem]{Remark}
\newcommand{\N}{\mathbb N}\newcommand{\Z}{\mathbb Z}
\newcommand{\Q}{\mathbb Q}\newcommand{\R}{\mathbb R}
\newcommand{\ipra}{\mathsf{iPRA}}
\newcommand{\PR}{\mathsf{PR}}
\newcommand{\F}{\mathcal F}
\newcommand{\E}{\mathcal E}
\newcommand{\diag}{\operatorname{diag}}
\newcommand{\tr}{\operatorname{tr}}
\newcommand{\norm}[1]{\left\lVert#1\right\rVert_F}
\newcommand{\eps}{\varepsilon}
\newcommand{\bound}{\operatorname{FB}}
\newcommand{\code}[1]{\texttt{\detokenize{#1}}}
\newcolumntype{P}[1]{>{\raggedright\arraybackslash}p{#1}}
\newcommand{\repo}[1]{\url{https://github.com/VladimirReshetnikov/ProveIt/blob/c39c388d4f87220ab9e758ae457520bbc0d74341/Logic/PrimitiveRecursiveArithmetic/#1}}
\lstset{basicstyle=\ttfamily\footnotesize,breaklines=true,columns=fullflexible,keepspaces=true,frame=single,rulecolor=\color{black!20},backgroundcolor=\color{black!2},showstringspaces=false}
\hypersetup{pdftitle={Stable observables from finite witnesses: an iPRA spectral framework},pdfauthor={Research prepared for Vladimir Reshetnikov},pdfsubject={Constructive arithmetic, rational spectral certificates, Lean formalization plan}}
\begin{document}
\hypersetup{pageanchor=false}
\begin{titlepage}
\vspace*{16mm}
{\Large\color{accent} CONSTRUCTIVE MATHEMATICS / ARITHMETIC FOUNDATIONS\par}
\vspace{12mm}
{\huge\bfseries\color{ink} Stable observables\par from finite witnesses\par}
\vspace{7mm}
{\Large An iPRA framework with rational spectral constructions\par}
\vspace{12mm}
{\large Prepared for Vladimir Reshetnikov\par}
{\large 15 November 2026\par}
\vspace{13mm}
\begin{minipage}{0.94\textwidth}
\textbf{Main results.} A uniform rational approximate spectral theorem; ordered eigenvalue names retaining multiplicities through collisions; a constructive principal matrix square root; and spectral projectors at a certified gap. Every construction has explicit finite data and primitive recursive control. The report also proves a precise limitation on complete numerical-bound certificate systems.

\textbf{Verification status.} The mathematical arguments are developed in full below. Five generic inference combinators have been compiled against the actual embedded \code{Prf PRAAx} calculus. The matrix-specific object-language proofs are specified and planned, but are not claimed as completed Lean formalizations. Exact rational reference programs and their verification records accompany the report.

\textbf{Source baseline.} ProveIt revision
\code{c39c388d4f87220ab9e758ae457520bbc0d74341}.
This is a targeted extension of the existing frameworks, not an independent verification of all 61 source packets or a claim of historical priority for spectral mathematics.
\end{minipage}
\vfill
{\small\color{accent} Complete source, executable artifacts, provenance, and a reproducible Lean subproject are included in the archive.}
\end{titlepage}
\hypersetup{pageanchor=true}
\pagenumbering{roman}
\begin{abstract}
Finite mathematical witnesses can be nonunique and unstable even when the quantity they determine is canonical. We organize embeddings into intuitionistic primitive recursive arithmetic around finite certificate producers and comparison estimates for the observations extracted from any two certificates. The abstract composition rule refines mechanisms already present in ProveIt's finite-contract, aggregate, and coherent-presentation packets. Its substantive new application here is a complete finite spectral construction for rational symmetric matrices, uniform in dimension.

A two-branch rational Jacobi step reduces an explicit off-diagonal energy without radicals. An elementary fuel bound yields exact rational orthogonal approximate diagonalizations. Rational subspace intersection then proves a perturbation estimate for their sorted diagonals, independent of the chosen bases. This produces extensional real eigenvalue names with multiplicities through collisions. Further finite comparison arguments give principal positive square roots and spectral projectors under a numerical gap, including the qualifications needed for represented positive inputs. The proof architecture exposes the precise quantifier-free invariants and the route to \code{Prf PRAAx}. Finally, a diagonal argument rules out a sound, complete effective system of numerical height-bound certificates for arbitrary primitive recursive fibers of cardinality at most one. General arithmetic difficulties that require new global bounds remain explicitly open.
\end{abstract}
\tableofcontents
\newpage
\pagenumbering{arabic}
\section{What this report adds}
\label{sec:scope}

The repository's own index distinguishes the kernel-checked ConciseIPRA development from the delivered proof packets, whose mathematical and Lean claims were not rerun during intake \cite{repoindex,packetindex}. At the pinned revision there are 61 packets, including ten foundations packets. Accordingly, a useful next framework should build on the existing finite-observation methods and supply a substantive missing mathematical comparison theorem.

The focus here is the passage from \emph{independent finite witnesses} to \emph{coherent, representation-independent observations}. A polynomial factorization, an approximate minimizer, or an approximate eigenbasis may change drastically when the input or precision changes. One should ask which information extracted from it has a quantitative comparison law. If such a law is proved for every pair of valid certificates, separate finite computations already suffice to construct a single coherent name for that information. Compatibility of the entire witnesses is unnecessary.

This principle is not presented as an unprecedented logical idea. The earlier aggregate packet obtains coherent extremal values from changing approximate minimizers; the coherent-methods packet derives coherence from uniqueness in finite presentations \cite{aggregates,coherent}. We give an explicit general interface and develop a new, multiplicity-sensitive spectral application with independent finite proofs of the required estimates.

\subsection{The principal conclusions}

All numerical error bounds in the following list are rational. Norm notation abbreviates squared rational inequalities, as defined in Section~\ref{sec:linear}.
\begin{enumerate}
\item Given a rational symmetric $n\times n$ matrix $A$ and $\eps>0$, a fixed primitive recursive algorithm returns a rational orthogonal matrix $Q$ and an increasing rational list $d$ such that
\[
 \norm{A-Q\diag(d)Q^{\mathsf T}}\le\eps.
\]
There is no dimension-dependent appeal to a separate existence theorem.
\item If two such certificates have residuals $\eta,\theta$ for inputs $A,B$, and $\norm{A-B}\le\delta$, their ordered observations satisfy
\[
 \max_{i<n}|d_i-e_i|\le\delta+\eta+\theta.
\]
This comparison gives regular primitive recursive eigenvalue names, with multiplicity, for represented symmetric matrices. It requires no gap between adjacent eigenvalues.
\item Positive rational matrices admit rational positive square-root approximants. The comparison estimate
\[
 \F(S-T)^2\le n\F(S^2-T^2),\qquad
 \F(X)=\sum_{i,j}X_{ij}^2,
\]
first holds for the explicitly constructed spectral certificates. A further finite approximation argument proves uniqueness against arbitrary positive represented roots. The representation correction for rational samples that are slightly negative is included.
\item If a threshold is separated from the spectrum by a supplied positive numerical margin, finite rational spectral projectors are coherent. Collisions within either selected block cause no problem. Projector stability is proved by a mixed-basis entrywise identity.
\item A general complete-certificate principle for numerical bounds is impossible even for primitive recursive fibers containing at most one element. The proof clarifies why finite checking cannot replace the missing global height-bound producers in other packets.
\end{enumerate}

\subsection{Comparison with the existing packets}

\begin{longtable}{@{}P{0.23\textwidth}P{0.34\textwidth}P{0.37\textwidth}@{}}
\toprule
Earlier development & Existing contribution & Additional result or retained boundary \\
\midrule\endhead
Finite contracts and finite observation \cite{contracts,observations} & Booleanized finite assertions, precision transformers, finite replay, and numerical presentations. & Reused as the arithmetical foundation. The new content is the comparison of independently generated spectral witnesses. \\
Finite aggregates and root supports \cite{aggregates,rootsupport} & Hausdorff root supports through collisions, spectral-radius names, and stable aggregate values. & Sorted symmetric spectra retain multiplicities and satisfy an explicit linear perturbation estimate. General complex root-multiset matching remains outside the result. \\
Uniform factorization \cite{factorization} & A finite producer of complete factorization data. & Existence of a list is separated from a comparison theorem between lists. The spectral theorem provides both in its symmetric subclass. \\
Coherent finite methods \cite{coherent} & Restriction-compatible finite data, uniqueness-based coherence, filtered lifting. & Approximate bases need neither be unique nor form a compatible tower; their ordered observations are compared directly. \\
Local-to-global methods \cite{localglobal} & A variable-length object-iPRA iteration rule. & The same induction interface is specialized to an energy-controlled computation and two-certificate comparison. \\
Euler--Rodrigues \cite{rotations} & Approximate axes without coherent axis selection, exact extraction with a gap, and quaternion data. & Approximate axes themselves are not claimed as new. We obtain an arbitrary-dimensional matrix framework and intrinsic spectral operations. \\
\bottomrule
\end{longtable}

The distinction between root \emph{supports} and multiplicity-bearing lists is essential. The lists $(0,0,1)$ and $(0,1,1)$ have the same support but different multiplicities; Hausdorff support distance cannot distinguish them. The earlier root-support report explicitly leaves this problem open \cite{rootsupport}. Our ordered comparison resolves the issue for spectra of real symmetric matrices. It is not a solution of the general complex-polynomial matching problem.

\subsection{Four levels of assertion}

We keep four levels separate throughout. An executable reference program demonstrates finite computations. A human arithmetic proof justifies a uniform theorem using defining equations and quantifier-free induction. A host Lean theorem establishes a statement in Lean's ambient mathematics. An embedded derivation has the actual type \code{Prf PRAAx Gamma phi}. Only the last is input to the repository's native transfer theorem.

The main mathematical arguments in this report are human arithmetic proofs. The archive contains compiled embedded inference combinators, but not embedded proofs of the spectral estimates. Section~\ref{sec:lean} specifies the latter's definitions, dependencies, induction invariants, and precise completion steps. Experimental checks are reported separately and are not used as substitutes for universal proofs.

\section{Arithmetic presentations and stable observations}
\label{sec:framework}

\subsection{The base theory and the status of infinite objects}

The base theory is first-order intuitionistic logic with equality, a function symbol for every fixed primitive recursive definition, its defining equations, the successor axioms, and induction for quantifier-free formulas. This is the repository's \code{PRAAx} presentation \cite{prsource}. We use the ordinary numerical domain. Rational numbers, finite arrays, matrices and certificates are numerical codes.

Signed integers may be coded by a sign and a magnitude. A rational is a signed numerator and a positive denominator; reduction to lowest terms is optional. Invalid inputs receive fixed default values, making every primitive recursive operation total. On valid codes, addition, multiplication, order, equality, and guarded reciprocal are the familiar cross-multiplication operations. Division at zero is assigned a default and is never used in the mathematical branch requiring a genuine reciprocal.

Finite lists can be coded by a length and iterated pairing. An array records its dimensions and row-major list. All loops over array entries have bounds computed from these dimensions. Bounded least search returns a designated default if no candidate passes. Sorting uses insertion or selection with rational comparison and a fixed tie break. These are direct primitive recursive definitions.

There is no numerical primitive recursive evaluator for all primitive recursive program codes. Every construction below is one fixed program, or a metatheoretic schema for a fixed family of input programs. For example, a represented matrix may use a fixed primitive recursive symbol $a(x,r,i,j)$, with numerical parameters $x,r,i,j$. We do not quantify inside arithmetic over arbitrary real-valued functions or assume that a code for an arbitrary PR definition can be evaluated by one universal PR function.

\begin{lemma}[Bounded formulas have Boolean presentations]
\label{lem:boolean}
For each fixed formula made from primitive recursive atomic predicates, propositional connectives, and bounded numerical quantifiers, one can construct a primitive recursive Boolean function $b$ and an $\ipra$ proof that the formula is equivalent to $b=1$.
\end{lemma}
\begin{proof}
Natural equality has a primitive recursive zero-test, and the defining equations and quantifier-free induction prove its correctness. Rational comparisons reduce to integer comparisons. Use Boolean multiplication for conjunction, maximum for disjunction, and Boolean negation for negation. Since the component predicates are decidable, implication is represented by the Boolean expression $\neg p\lor q$; this equivalence is intuitionistically valid for a decidable antecedent.

For a bounded universal quantifier use the product of the Boolean values over the stated finite range. For a bounded existential use their maximum, or the flag returned by a bounded scan. For the universal fold, fix the final bound as a parameter. Prove that a fold equal to one forces each indexed bit to equal one by quantifier-free prefix induction with that index free. Conversely, work under the fixed hypothesis that all bits below the final bound equal one, and prove by quantifier-free prefix induction that the partial product is one. Generalize the free index only after the induction. The existential case uses a bounded scan together with its stored witness or its all-zero flag. Thus the proof does not induct on the quantified formula whose Boolean presentation is being established. Nested bounded quantifiers are handled by composing the folds, not by adding an induction scheme for unbounded formulas.
\end{proof}

Consequently assertions such as ``the matrix is symmetric'', ``the rational array is sorted'', and ``every recorded transition obeys its update rule'' are admissible induction predicates after compilation. The fact that a paper writes a bounded assertion with $\forall i<n$ does not add logical strength. By contrast, ``all future precisions work'' or ``all rational points have already been found'' cannot simply be inserted into the invariant.

\subsection{A certificate interface}

\begin{definition}[Finite presentation with observable comparison]
\label{def:interface}
Fix primitive recursive numerical functions and predicates:
\begin{enumerate}
\item an input-validity predicate $V(x)$;
\item a finite certificate predicate $C(x,r,w)$;
\item a producer $P(x,r)$ and observation $o(x,r,w)$;
\item a finite input-comparison predicate $D(x,y,d)$ and an output-comparison predicate $R(u,v,e)$;
\item a primitive recursive rational error budget $\omega(d,r,s)$.
\end{enumerate}
The interface includes actual arithmetic proofs of
\begin{align}
 V(x)&\longrightarrow C(x,r,P(x,r)), \label{eq:producer}\\
 D(x,y,d)\wedge C(x,r,w)\wedge C(y,s,z)
 &\longrightarrow R(o(x,r,w),o(y,s,z),\omega(d,r,s)).
 \label{eq:comparison}
\end{align}
All displayed quantifiers over free numerical variables may be closed universally after the finite argument. The comparison concerns arbitrary accepted certificates, not merely one favored execution.
\end{definition}

\begin{theorem}[Stable-observation composition]
\label{thm:observable}
Under the two proved interface clauses, define $O(x,r)=o(x,r,P(x,r))$. Then $\ipra$ proves the same comparison for $O(x,r),O(y,s)$ whenever the two inputs are valid and satisfy $D(x,y,d)$.
\end{theorem}
\begin{proof}
Instantiate producer correctness separately at $(x,r)$ and $(y,s)$, substitute the two resulting certificates in the comparison theorem, and apply implication elimination twice. Every witness is an explicit primitive recursive term. No choice, infinite compatibility assertion, or new induction is involved.
\end{proof}

The proof is deliberately elementary. The difficulty of using the theorem is entirely in proving the producer and the comparison estimates. Same-input comparison establishes independence from witness selection for that input. Cross-input comparison is additionally required for perturbation bounds and independence from the chosen name of the input. It is a material error to infer the second from the first without proof.

For a rational metric presentation a typical comparison is
\[
 d_Y(O(x,r),O(y,s))\le Ld_X(x,y)+e_r+e_s,
 \qquad e_r=2^{-r}.
\]
For square roots the comparison is naturally expressed by squared errors instead of a Lipschitz constant. A precision transformer then chooses certificate tolerances sufficiently small to meet the requested observation tolerance. The framework does not require a uniform exponent or one metric for every application.

\subsection{Regular names and finite meaning}

A regular rational matrix name is a fixed primitive recursive family $A_r$ with the arithmetic contract
\[
 \bound(A_r-A_s,e_r+e_s)\qquad(r,s\in\N),
 \qquad e_r=2^{-r}.
\]
Here $\bound(X,a)$ means $a\ge0$ and $\F(X)\le a^2$. For scalar or vector observations use the analogous absolute-value or coordinatewise comparison. The represented object is a useful semantic interpretation of this family; the internal theorem consists of the finite rational inequalities for arbitrary numerical $r,s$.

Two regular names represent the same matrix when
\[
 \bound(A_r-B_s,e_r+e_s)\qquad\text{for all }r,s.
\]
This is a formula with unbounded universal quantifiers, not a decidable test on names. Its transitivity follows by comparing through a common arbitrarily fine observation. There is a fully arithmetic justification: if the proposed rational squared bound fails, the positive rational gap determines a concrete later precision for which the intermediate error is too small to allow the failure. Thus no unrestricted limiting principle is needed.

We will use this finite-slack argument repeatedly. For a positive rational $a=p/q$, with $p,q\ge1$, one can ensure $2^{-k}<a$ by the deliberately excessive primitive recursive choice $k=q+1$, since $2^{q+1}>q$. More efficient binary-length choices are optional. Whenever a finite bound contains a rational coefficient $K$, choose a precision with $K2^{-k}$ below the prescribed positive rational gap. The witness is numerical and primitive recursive.

\subsection{Explicit fuel and quantifier-free invariants}

Let $s_0(x)$ and $\operatorname{step}(x,s)$ be fixed primitive recursive functions. Define their iteration by primitive recursion in $t$. A proof of a finite computation should exhibit a Boolean invariant $I(x,t,s_t)$, prove its base and step cases, and specialize quantifier-free induction at a primitive recursive fuel $T(x)$.

For the spectral computation the decisive numerical part is
\[
 \E(A_t)\le\eps^2
 \quad\lor\quad
 m\E(A_t)+t\eps^2\le m\F(A),
 \qquad m=n(n-1).
\]
All values are rational finite sums. Orthogonality, symmetry and the relation between the current matrix and the initial matrix join this invariant by finite Boolean conjunction. Once the left alternative holds the transition freezes, so the invariant remains valid without recording a least stopping time.

An ambient statement that an algorithm terminates does not supply this object-theory proof. Similarly, a certificate checker and a semantic proof of its soundness do not supply a total producer. These distinctions are part of the interface, not optional documentation. Section~\ref{sec:certificate-boundary} shows a setting where no general complete certificate producer can exist.

\subsection{Composition across domains}

Suppose one stage produces an observation used as the input to a second stage. To compose them, request the first stage's output at the tolerance required by the second stage's input-comparison theorem, then add or otherwise combine their error budgets using that theorem. The requested precision must be an explicit primitive recursive function of the final demand and the finite input bounds.

This applies to algebraic identities, norms, finite geometry, matrix functions and analytic estimates equally. What varies is the object being observed and the modulus. For example: a factorization can yield a root support; a minimization can yield an infimum value; a diagonalization can yield ordered eigenvalues; a noncanonical basis for a separated spectral block can yield a canonical projector. The interface organizes these differences instead of asserting that every mathematical object has a computable canonical representative.

\section{Finite rational linear algebra}
\label{sec:linear}

\subsection{Norm notation without irrational operations}

For a rational matrix put
\[
 \F(A)=\sum_{i,j<n} A_{ij}^2,
 \qquad \E(A)=\sum_{\substack{i,j<n\\i\ne j}}A_{ij}^2.
\]
Write $\bound(A,a)$ for $a\ge0$ and $\F(A)\le a^2$. The notation $\norm A\le a$ means precisely $\bound(A,a)$. No primitive function for an irrational norm is required. A displayed comparison involving norms can always be read as the rational squared assertion stated alongside it. Likewise $\|v\|^2=\sum_i v_i^2$ is rational.

The relevant finite inequalities have elementary sum-of-squares proofs. For rational lists $u,v$ of equal length,
\[
 \left(\sum_i u_i^2\right)\left(\sum_i v_i^2\right)
 -\left(\sum_i u_iv_i\right)^2
 =\sum_{i<j}(u_iv_j-u_jv_i)^2\ge0.
\]
Expanding and collecting terms proves the identity by finite distributivity. Prefix induction verifies the expansion uniformly in length. Taking signs or absolute values uses only decidable rational comparisons. This proves finite Cauchy--Schwarz, the triangle inequality for Frobenius bounds, and the submultiplicative estimate
\[
 \F(AB)\le\F(A)\F(B).
\]
For the last estimate apply Cauchy--Schwarz to each matrix entry and sum over the row and column indices.

\begin{lemma}[Quadratic-form bound]
\label{lem:qform}
If $\bound(R,a)$, then for every rational vector $v$,
\[
 |v^{\mathsf T}Rv|\le a\sum_i v_i^2.
\]
\end{lemma}
\begin{proof}
Apply the finite inequality to the two arrays $R_{ij}$ and $v_iv_j$. Their second squared sum is
\[
 \sum_{i,j}v_i^2v_j^2=\left(\sum_i v_i^2\right)^2.
\]
The squared asserted inequality follows. Both sides of the unsquared comparison are nonnegative rationals; if the comparison failed, their squares would have the same strict ordering, a contradiction.
\end{proof}

A rational orthogonal certificate will include both identities
$Q^{\mathsf T}Q=I$ and $QQ^{\mathsf T}=I$. The algorithm constructs $Q$ as a product of explicit plane rotations, so both identities are immediate invariants. Including both in a checker also avoids requiring a separate invertibility argument there. Finite expansion gives
\[
 \F(Q^{\mathsf T}AR)=\F(A)
\]
whenever $Q,R$ are square orthogonal matrices of the appropriate dimensions. Indeed sum first over the outer index and use $\sum_i Q_{ki}Q_{\ell i}=\delta_{k\ell}$, and then do the same on the right.

\subsection{Constructive intersection by rational elimination}

\begin{lemma}[Finite intersection]
\label{lem:intersection}
Let $Q,R$ be rational orthogonal $n\times n$ matrices and $0\le k<n$. The span of the first $k+1$ columns of $Q$ and the span of the last $n-k$ columns of $R$ contain a common nonzero rational vector. A fixed primitive recursive procedure constructs it.
\end{lemma}
\begin{proof}
Form the $n\times(n+1)$ rational matrix
\[
 W=[\,Q_{\{0,\ldots,k\}}\mid -R_{\{k,\ldots,n-1\}}\,].
\]
Perform Gaussian elimination, scanning the remaining rows and columns in a fixed order for a nonzero rational pivot. This is a bounded scan; division occurs only after a successful nonzero test. Each pivot eliminates one column below its row. At most $n$ pivot columns can occur, so one of the $n+1$ columns is free. Set its free coordinate to one, all other free coordinates to zero, and back substitute to obtain a nonzero rational nullvector $(u,v)$.

The nullvector equation says $Q_{\le k}u=R_{\ge k}v=:x$. If $x=0$, multiplying by the transpose of the respective column family gives $u=0$ and $v=0$, because its columns are orthonormal. This contradicts the selected free coordinate. Thus $x\ne0$.

The procedure performs at most $n$ pivot stages, each containing only loops bounded by the array dimensions, followed by at most $n$ back-substitution stages. Its row-operation identity, pivot-position constraints and nullvector equations are finite predicates. Their Booleanized versions are quantifier-free invariants. No decision of the rank of a real matrix is made.
\end{proof}

\subsection{Positive rational matrices have finite certificates}

\begin{definition}[Positive certificate]
For a rational symmetric $A$, a positive certificate is rational data $(L,a)$ with
\[
 A=L\diag(a)L^{\mathsf T},\qquad a_i\ge0\ (i<n).
\]
The factor $L$ need not be orthogonal or triangular in the certificate format. It may include row permutations from elimination. The displayed identities are checked exactly.
\end{definition}

Such a certificate implies $v^{\mathsf T}Av\ge0$ by a finite sum of nonnegative squares. Conversely, finite rational elimination can either construct a certificate or construct a counterexample.

\begin{theorem}[Finite positivity alternative]
\label{thm:ldl}
A primitive recursive algorithm, uniform in $n$, returns for every rational symmetric $A$ either a positive certificate or a rational vector $v$ with $v^{\mathsf T}Av<0$. Hence rational positive semidefiniteness has a complete finite certificate interface.
\end{theorem}
\begin{proof}
For dimension zero return the empty certificate. If some diagonal entry is negative, the corresponding coordinate vector is a counterexample. If every diagonal entry is zero, then either every entry is zero, yielding the zero diagonal certificate, or some off-diagonal $b=A_{ij}$ is nonzero. In the latter case
\[
 v=e_i-\operatorname{sgn}(b)e_j
 \quad\text{satisfies}\quad v^{\mathsf T}Av=-2|b|<0.
\]
All cases are decided by finite scans.

In the remaining case choose a positive diagonal entry, move it to the first position by a permutation, and write
\[
 A=\begin{pmatrix}a&v^{\mathsf T}\\v&B\end{pmatrix},\qquad a>0.
\]
Compute the rational Schur complement $C=B-vv^{\mathsf T}/a$. Repeat the same process on $C$. If the recursive result is a negative vector $y$, lift it to
\[
 x=\begin{pmatrix}-v^{\mathsf T}y/a\\y\end{pmatrix};
 \qquad x^{\mathsf T}Ax=y^{\mathsf T}Cy<0.
\]
If $C=K\diag(c)K^{\mathsf T}$ with $c\ge0$, use
\[
 L=\begin{pmatrix}1&0\\v/a&K\end{pmatrix},\qquad
 A=L\diag(a,c)L^{\mathsf T}.
\]
Undo the permutation. This proves correctness of the elimination step.

There are at most $n$ dimension-reduction stages. To implement it as ordinary primitive recursion, retain a stack of the pivots, columns and permutations; then reconstruct the certificate or negative vector in at most $n$ reverse stages. The invariant records the exact block congruence with the current Schur matrix and the positivity of the processed pivots. These are bounded rational assertions. The final alternative, including the produced negative vector in the failure branch, is likewise finite. Thus the argument can be expressed by quantifier-free induction over the two bounded passes. This avoids an induction hypothesis quantifying over all matrices of the current size.
\end{proof}

This theorem concerns exact rational inputs. It does not decide positivity of an arbitrary represented real matrix. Section~\ref{sec:psdnames} gives the separate correction needed for positive real names with slightly indefinite rational samples.

\section{A rational approximate spectral theorem}
\label{sec:jacobi}

The classical Jacobi method motivates the construction \cite{jacobinotes}, but none of its convergence or spectral theorems is assumed. We prove a finite, entirely rational variant with a deliberately elementary fuel bound.

\subsection{An exact energy identity}

Fix a symmetric rational matrix $A$ and indices $i<j$. Its pivot block is
\[
 \begin{pmatrix}a&b\\b&d\end{pmatrix},\qquad h=a-d.
\]
Let $G$ be the identity outside these coordinates, with pivot block
\[
 G_{ij}=\begin{pmatrix}c&-s\\s&c\end{pmatrix},
 \qquad c^2+s^2=1.
\]
Set $A'=G^{\mathsf T}AG$. Direct multiplication gives
\begin{equation}
 b'=(c^2-s^2)b-csh. \label{eq:pivot}
\end{equation}
For each $k$ outside the pivot, the pair $(A_{ki},A_{kj})$ is multiplied by an orthogonal $2\times2$ matrix and preserves its squared sum. Symmetry supplies the matching row pair. All remaining off-diagonal entries are unchanged. Therefore
\begin{equation}
 \E(A')=\E(A)-2b^2+2(b')^2. \label{eq:energy}
\end{equation}
This identity isolates the entire progress argument in two rational numbers $h,b$.

\subsection{Two rational branches}

\begin{lemma}[Rational pivot improvement]
\label{lem:pivot}
For any rational $b\ne0$ and $h$, the following decidable two-branch construction produces rational $c,s$ with $c^2+s^2=1$ and $(b')^2\le b^2/2$. Consequently $\E(A')\le\E(A)-b^2$.
\end{lemma}
\begin{proof}
If $h^2\ge b^2$, then $h\ne0$. Put
\[
 t=\frac b{4h},\qquad
 c=\frac{1-t^2}{1+t^2},\qquad s=\frac{2t}{1+t^2}.
\]
The denominator $1+t^2$ is positive, and expansion gives $c^2+s^2=1$. With $u=t^2$ we have $0\le u\le1/16$. Substitution in \eqref{eq:pivot} gives
\[
 b'=b\frac{1-11u+2u^2}{2(1+u)^2}.
\]
The numerator is nonnegative because $1-11u\ge5/16$. Also
\[
 1-11u+2u^2\le(1+u)^2
 \quad\Longleftrightarrow\quad u^2\le13u,
\]
which holds on this interval. The factor multiplying $b$ lies between zero and $1/2$, so $(b')^2\le b^2/4$.

If $h^2<b^2$, take $c=3/5$ and compare $s=4/5$ with $s=-4/5$. Choose the one giving the smaller rational square of the transformed pivot. The two pivots are
\[
 b'_+=\frac{-7b-12h}{25},\qquad
 b'_- =\frac{-7b+12h}{25}.
\]
Their mean square equals
\[
 \frac{(b'_+)^2+(b'_-)^2}{2}
 =\frac{49b^2+144h^2}{625}
 <\frac{193}{625}b^2<\frac12b^2.
\]
The smaller square is at most the mean. This proves the common bound in both branches. Equation~\eqref{eq:energy} gives the claimed decrease.
\end{proof}

No irrational angle, square root, or real equality test occurs. The high-difference branch guards its division; the low-difference branch has fixed positive denominators. On the unused input $b=0$ the totalized step returns the identity rotation.

\subsection{The finite algorithm and its proof}

\begin{definition}[Diagonalization certificate]
\label{def:diagonalcertificate}
For symmetric rational $A$ and rational $\eps>0$, a certificate consists of a rational orthogonal $Q$ and an increasing rational list $d$ such that, with $D=\diag(d)$,
\[
 B=Q^{\mathsf T}AQ,\qquad d_i=B_{ii},\qquad
 \F(B-D)\le\eps^2.
\]
In particular the residual has zero diagonal. Both orthogonality identities are part of the finite format.
\end{definition}

\begin{theorem}[Uniform rational approximate diagonalization]
\label{thm:jacobi}
There is a fixed primitive recursive producer of diagonalization certificates. For $n\ge2$ it needs at most
\begin{equation}
 T=n(n-1)\left(\left\lceil\frac{\F(A)}{\eps^2}\right\rceil+1\right)
 \label{eq:fuel}
\end{equation}
rotation stages. Dimensions zero and one need no rotations.
\end{theorem}
\begin{proof}
Set $A_0=A,Q_0=I$ and $m=n(n-1)$. At each stage, freeze if $\E(A_t)\le\eps^2$. Otherwise choose the first pair $i<j$ maximizing $A_{t,ij}^2$, apply Lemma~\ref{lem:pivot}, and set
\[
 A_{t+1}=G_t^{\mathsf T}A_tG_t,
 \qquad Q_{t+1}=Q_tG_t.
\]
The maximum exists by a finite scan. At an active stage it is nonzero. There are $m/2$ unordered pairs and energy counts each twice; hence
\[
 b^2\ge\E(A_t)/m>\eps^2/m.
\]
Every active stage reduces energy by at least $\eps^2/m$.

The exact invariant includes
\[
 Q_t^{\mathsf T}Q_t=Q_tQ_t^{\mathsf T}=I,\quad
 A_t=Q_t^{\mathsf T}AQ_t,\quad A_t=A_t^{\mathsf T},\quad
 \F(A_t)=\F(A).
\]
The progress component is the quantifier-free rational disjunction
\begin{equation}
 \E(A_t)\le\eps^2
 \quad\lor\quad
 m\E(A_t)+t\eps^2\le m\F(A). \label{eq:invariant}
\end{equation}
It holds at zero. A frozen stage preserves its first alternative. At an active stage use the decrease estimate to prove the second alternative at $t+1$ if the first does not already hold. This proves the invariant by quantifier-free induction.

Write $K=\lceil\F(A)/\eps^2\rceil+1$, so $K\eps^2>\F(A)$. At $T=mK$ the second alternative in \eqref{eq:invariant}, together with nonnegative energy, would imply $K\eps^2\le\F(A)$, a contradiction. Thus the desired residual inequality holds.

Finally sort the rational diagonal entries, breaking ties by the original index. Apply the corresponding permutation simultaneously to the columns of $Q_T$ and the rows and columns of $A_T$. Exact orthogonality, the transformed-matrix identity, and off-diagonal energy are preserved. This yields the stated certificate. If $n=0$ or $1$, use the identity matrix and the already sorted diagonal.

Every update, scan, comparison, permutation and fuel computation is primitive recursive on its finite code. Booleanize the finite matrix equalities and combine them with \eqref{eq:invariant}. This identifies the actual admissible induction formula, rather than merely asserting that some termination proof exists.
\end{proof}

One can stop the implementation as soon as its state freezes. This is exactly equivalent to the specified $T$-stage primitive recursion, because all subsequent states are unchanged. The reference program uses this optimization.

\subsection{A grid refinement controlling rational denominators}
\label{sec:grid}

The preceding proof is simple, but repeated exact ratios can create very large denominators. The following variant keeps each rotation denominator bounded in terms of the original finite input and precision.

Put
\[
 M=1+\lceil\F(A)\rceil,
 \qquad L=\max\left(1,\left\lceil\frac{128Mn}{\eps}\right\rceil\right).
\]
Then $\norm{A_t}\le M$, so $|b|\le M$ and $|h|\le2M$. In the high-difference branch replace $t=b/(4h)$ by its truncation $v$ toward zero on the $1/L$ grid. Thus $|v|\le|t|\le1/4$ and $|v-t|<1/L$. Use the same rational circle parametrization at $v$. The low-difference branch is unchanged.

\begin{proposition}[Finite-grid progress]
\label{prop:grid}
The grid variant decreases active-stage energy by at least $b^2/2$ and therefore succeeds with twice the fuel in \eqref{eq:fuel}. Each rotation entry has a denominator at most $\max(5,2L^2)$ before reduction.
\end{proposition}
\begin{proof}
For $|u|,|v|\le1$ the exact rational identities
\[
 c(u)-c(v)=\frac{2(v^2-u^2)}{(1+u^2)(1+v^2)},\qquad
 s(u)-s(v)=\frac{2(u-v)(1-uv)}{(1+u^2)(1+v^2)}
\]
give $|c(u)-c(v)|,|s(u)-s(v)|\le4|u-v|$. Since $|c|,|s|\le1$, the transformed-pivot expression changes by at most
\[
 (16|b|+8|h|)|u-v|\le32M|u-v|.
\]
At an active maximal pivot, $|b|>\eps/n$, because $b^2>\eps^2/[n(n-1)]$. The rounding changes the ideal pivot by at most
\[
 32M/L\le\eps/(4n)<|b|/4.
\]
Its new absolute value is therefore at most $3|b|/4$. The low branch has squared pivot at most $b^2/2$, which is also at most $9b^2/16$. Equation~\eqref{eq:energy} now gives a decrease of at least $7b^2/8$, and hence of at least $b^2/2$. Use the invariant with $2m$ in place of $m$ and double the fuel.

Writing $v=j/L$, the rotation entries have common denominator $L^2+j^2\le2L^2$. In the low branch the denominator is $5$. This proves the denominator claim.
\end{proof}

If $H=\max(5,2L^2)$, the product $Q_t$ has some common denominator at most $H^t$; this need not mean that its denominator divides the particular integer $H^t$. Products of the individual rotation denominators provide the claimed bound. For an input common denominator $D_0$, the transformed matrix has one at most $D_0H^{2t}$. Entry magnitudes remain bounded by $1$ in $Q_t$ and $M$ in $A_t$. Thus the entries admit reduced or common-denominator representatives whose numerator and denominator bit lengths have explicit bounds linear in $t\log H$ plus the input sizes. This is a storage bound when fractions are reduced, as in the reference program, or when a common denominator is explicitly maintained; it is not a bound on arbitrarily redundant unreduced intermediate encodings. We make no polynomial-time claim in binary precision: the elementary fuel itself grows like $\eps^{-2}$.

\section{Ordered spectra that remain coherent through collisions}
\label{sec:spectra}

\subsection{Comparison of arbitrary certificates}

\begin{theorem}[Finite spectral comparison]
\label{thm:weyl}
Let $A,B$ be rational symmetric $n\times n$ matrices. Suppose $(Q,d)$ and $(R,e)$ are diagonalization certificates with errors $\eta,\theta$, respectively, and $\bound(A-B,\delta)$. Then for every $k<n$,
\begin{equation}
 |d_k-e_k|\le\delta+\eta+\theta. \label{eq:weyl}
\end{equation}
No comparison between $Q$ and $R$ is assumed.
\end{theorem}
\begin{proof}
For $n=0$ there is no coordinate to check. Otherwise use Lemma~\ref{lem:intersection} to construct a nonzero rational vector $x$ in the span of the first $k+1$ columns of $Q$ and the last $n-k$ columns of $R$.

In the first coordinates the sortedness of $d$ gives
\[
 x^{\mathsf T}Q\diag(d)Q^{\mathsf T}x\le d_k x^{\mathsf T}x.
\]
The residual changes this quadratic form by at most $\eta x^{\mathsf T}x$, by Lemma~\ref{lem:qform}. Similarly the second coordinate representation gives
\[
 x^{\mathsf T}Ax\le(d_k+\eta)x^{\mathsf T}x,
 \qquad
 x^{\mathsf T}Bx\ge(e_k-\theta)x^{\mathsf T}x.
\]
Subtract and apply Lemma~\ref{lem:qform} to $B-A$:
\[
 (e_k-d_k-\eta-\theta)x^{\mathsf T}x
 \le x^{\mathsf T}(B-A)x\le\delta x^{\mathsf T}x.
\]
Because $x\ne0$, the rational $x^{\mathsf T}x$ is strictly positive. Divide by it to obtain $e_k-d_k\le\delta+\eta+\theta$. Interchange the two certificates and repeat the argument for the reverse inequality.

The proof used only bounded elimination, finite sums and rational order. The construction of $x$ is not an invocation of a real-vector-space dimension theorem with an unexamined witness principle.
\end{proof}

This is the main comparison theorem demanded by the framework. It is stronger than a statement that the diagonalizer happens to converge along one selected execution. It applies to independently produced certificates at different precisions, with different rotations and tie breaks, and to different rational approximants of the input.

\subsection{Regular eigenvalue names}

Suppose a symmetric represented matrix has fixed primitive recursive entry names $a_{ij}(r)$ satisfying
\[
 |a_{ij}(r)-a_{ij}(s)|\le2^{-r}+2^{-s},
\]
and satisfy the arithmetic symmetry contract $|a_{ij}(r)-a_{ji}(s)|\le2^{-r}+2^{-s}$ for every pair of indices and precisions. To simplify the finite contract one may require the supplied rational arrays to be symmetric. Otherwise replace a sample by its average with its transpose; the entrywise error budget is preserved.

For an output precision $r$, sample at $r+n+4$, symmetrize, and call the result $A[r]$. Since $n\le2^n$ for $n\ge1$, its pairwise Frobenius comparison has budget
\[
 \bound\bigl(A[r]-A[s],\,2^{-r-4}+2^{-s-4}\bigr).
\]
Run Theorem~\ref{thm:jacobi} with tolerance $2^{-r-4}$ and denote the sorted diagonal by $d[r]$. Theorem~\ref{thm:weyl} yields
\begin{equation}
 |d[r]_i-d[s]_i|\le2^{-r-3}+2^{-s-3}. \label{eq:name}
\end{equation}
Thus every coordinate is a regular primitive recursive real name with spare precision margin. Dimension zero gives the empty family. A matrix already given by a Frobenius-regular name needs only a fixed precision shift rather than the entrywise dimension shift.

The names are ordered, count multiplicity, and are extensional with respect to the input names. For the last assertion compare the two rational input samples in Theorem~\ref{thm:weyl}, use their equality contract, and request sufficiently fine finite indices. There is no need to synchronize the bases. For represented matrices a Frobenius perturbation of at most $\delta$ gives
\[
 |\lambda_i(A)-\lambda_i(B)|\le\delta
\]
in the finite-slack interpretation. The full arithmetic theorem is \eqref{eq:weyl} together with the explicit precision shifts, rather than an added axiom about real limits.

\subsection{Why these names are eigenvalues}

One must identify the stable observations mathematically, not merely call them eigenvalues. This follows from a finite determinant estimate.

\begin{lemma}[Coefficient residual]
\label{lem:coeff}
Suppose $\bound(A,M)$ with $M\ge1$ and $(Q,d)$ is a diagonalization certificate of error $\eps$. For $1\le j\le n$, the coefficient of $X^{n-j}$ in $\det(XI-A)$ differs from that in $\prod_i(X-d_i)$ by at most
\begin{equation}
 K_j(M,n)\eps,
 \qquad K_j(M,n)=\binom nj j!\,j M^{j-1}. \label{eq:coefficient}
\end{equation}
\end{lemma}
\begin{proof}
Let $C=Q^{\mathsf T}AQ$ and $D=\diag(d)$. Every entry of $C,D$ has absolute value at most $M$, because $\F(C)=\F(A)$ and $D$ retains only the diagonal. Every entry of $C-D$ has absolute value at most $\eps$.

The coefficient in question is a signed sum of the $\binom nj$ principal minors of size $j$. Expand each determinant into its $j!$ signed products. Compare a product for $C$ with its product for $D$ by replacing one factor at a time. Each replacement costs at most $\eps M^{j-1}$, and there are $j$ replacements. Sum these bounds over the products and minors. Finally
\[
 \det(XI-Q^{\mathsf T}AQ)=\det(XI-A)
\]
follows from determinant multiplication and both exact orthogonality identities. These determinant identities themselves have finite permutation-sum proofs; their variable-size sums are bounded folds. The leading coefficients are both one.
\end{proof}

For represented input, a uniform magnitude bound is available from the stage-zero samples: choose an integer
\[
 M\ge1+n\left(\max_{i,j}|a_{ij}(0)|+2\right).
\]
Every later sample then has Frobenius bound $M$. All coefficient constants in \eqref{eq:coefficient} are explicit primitive recursive rational data. Request a precision at which these constants times the input and diagonalization errors are below the desired coefficient tolerance. Polynomial coefficients in the diagonal observations also have finite product moduli, obtained by the same telescoping argument. We obtain equality of represented coefficient names:
\begin{equation}
 \det(XI-A)=\prod_{i<n}(X-\lambda_i(A)). \label{eq:factorization}
\end{equation}

This supplies $n$ ordered real root names and their coefficient factorization, retaining repetitions. Every complex characteristic root is real as well: positive apartness of its imaginary part would give the same positive lower bound for the modulus of every factor, making their finite product apart from zero. Decidable rational approximation comparisons then force its imaginary part to be zero. We do not assert the disjunctive statement that every separately supplied exact root comes with an index identifying it with one of the ordered names. For $X^2-t^2$, selecting a matching index for the root $t$ from the list $(-|t|,|t|)$ would already give the generally unavailable comparison $t\le0\lor t\ge0$. We did not assume a fundamental theorem of algebra or an exact orthogonal eigenbasis in obtaining it. The theorem concerns the factorization and stable eigenvalue names; it does not assert a primitive recursive choice of an exact eigenbasis for arbitrary real names.

\subsection{Multiplicity information and the remaining decision problem}

For
\[
 A(t)=\begin{pmatrix}t&0\\0&-t\end{pmatrix},
\]
the ordered spectral observations name $(-|t|,|t|)$. At $t=0$ the two entries coincide, and the family remains continuous. A rule that assigns the first eigenvector to the smaller eigenvalue must exchange coordinate axes as the sign of $t$ changes. The matrix of witnesses can therefore fail to be continuous even though the ordered observations are stable.

Retaining multiplicity means producing $n$ ordered names, including repeated values in the factorization. It does \emph{not} mean deciding whether two represented values are equal, deciding a rank, or computing a partition of the index set into equal-eigenvalue blocks. Those tasks require additional information. A supplied gap allows a finite block decision, developed in Section~\ref{sec:projectors}; no gap is needed for the names themselves.

\section{Principal positive square roots without coherent eigenvectors}
\label{sec:squareroots}

This application goes beyond extracting a list of eigenvalues. We reconstruct a matrix from noncanonical spectral witnesses and prove that the reconstructed matrices are coherent. A separate argument establishes uniqueness among all positive represented roots.

\subsection{A finite square-root certificate}

Let $A$ be rational positive semidefinite, with a certificate from Theorem~\ref{thm:ldl}, and let $\eta>0$. For $n\ge1$, diagonalize $A$ with error $\eta/2$. Every diagonal observation is nonnegative:
\[
 d_i=q_i^{\mathsf T}Aq_i\ge0,
\]
where $q_i$ is the corresponding rational column of $Q$. This uses the input positive certificate and exact orthogonality, not the spectral theorem.

Choose an integer $U\ge1+\norm A$, for example
$U=2+\lceil\F(A)\rceil$. This example is valid because $2+x\ge1+\sqrt{x}$ for every nonnegative rational $x$, as a squared arithmetic comparison. Put
\begin{equation}
 H=\left\lceil\frac{2n(2U+1)}{\eta}\right\rceil.
 \label{eq:sqrtgrid}
\end{equation}
Then $H\ge1$ and $0\le d_i\le U\le U^2$. By a bounded scan, or a bounded binary search, find the largest integer $j_i\le UH$ with
\[
 j_i^2\le d_iH^2.
\]
Define $s_i=j_i/H$. If $j_i=UH$, maximality and $d_i\le U^2$ imply $s_i^2=d_i$. Otherwise the next integer fails the comparison, giving
\[
 0\le d_i-s_i^2<\frac{2j_i+1}{H^2}
 \le\frac{2U+1}{H}\le\frac{\eta}{2n}.
\]
Set
\begin{equation}
 S=Q\diag(s)Q^{\mathsf T}. \label{eq:sqrtapprox}
\end{equation}
This matrix has an exact positive spectral certificate.

\begin{proposition}[Finite positive root approximation]
\label{prop:sqrtproducer}
The construction is primitive recursive and satisfies $\bound(S^2-A,\eta)$.
\end{proposition}
\begin{proof}
In $Q$ coordinates, $A=D+R$ with $R$ having zero diagonal and $\F(R)\le\eta^2/4$. The square of \eqref{eq:sqrtapprox} is $Q\diag(s_i^2)Q^{\mathsf T}$. Therefore
\[
 \F(S^2-A)
 =\sum_i(s_i^2-d_i)^2+\F(R)
 \le n\left(\frac{\eta}{2n}\right)^2+\frac{\eta^2}{4}
 \le\frac{\eta^2}{2}\le\eta^2.
\]
The sum has at least one term only when $n\ge1$; dimension zero is handled by the empty output. All searches have explicit integer bounds from \eqref{eq:sqrtgrid}. Binary search may be run for the bit length of $UH+1$ plus one and then frozen, or replaced by the literal finite scan.
\end{proof}

\subsection{Stability in mixed bases}

\begin{theorem}[Square-root certificate stability]
\label{thm:sqrtstable}
Let
\[
 S=Q\diag(s)Q^{\mathsf T},\qquad
 T=R\diag(t)R^{\mathsf T},
\]
where $Q,R$ are rational orthogonal matrices and $s_i,t_j\ge0$. Then the entirely rational inequality
\begin{equation}
 \F(S-T)^2\le n\F(S^2-T^2) \label{eq:sqrtstable}
\end{equation}
holds. In particular, if $\bound(S^2-T^2,\delta)$ and $n\ge1$, then
$\F(S-T)\le n\delta$.
\end{theorem}
\begin{proof}
Put $V=Q^{\mathsf T}R$. In the mixed bases,
\begin{align*}
 (Q^{\mathsf T}(S-T)R)_{ij}&=(s_i-t_j)V_{ij},\\
 (Q^{\mathsf T}(S^2-T^2)R)_{ij}&=(s_i^2-t_j^2)V_{ij}.
\end{align*}
For nonnegative rationals $s,t$, $|s-t|\le s+t$, so
$(s-t)^2\le|s^2-t^2|$. Consequently
\[
 \F(S-T)
 \le\sum_{i,j}|(s_i^2-t_j^2)V_{ij}|\,|V_{ij}|.
\]
Apply finite Cauchy--Schwarz to the two arrays on the right and use
$\sum_{i,j}V_{ij}^2=n$. Orthogonal invariance proves \eqref{eq:sqrtstable}. If $\F(S^2-T^2)\le\delta^2$, then
$\F(S-T)^2\le n\delta^2\le n^2\delta^2$ for $n\ge1$. Nonnegative rational square comparison gives the convenient weaker bound $\F(S-T)\le n\delta$.
\end{proof}

For fixed rational $A$, construct $S_r$ with residual tolerance
\[
 \eta_r=\frac{2^{-2r-4}}{n+1}.
\]
The squares satisfy $\bound(S_r^2-S_s^2,\eta_r+\eta_s)$. The theorem gives
\[
 \F(S_r-S_s)\le n(\eta_r+\eta_s)
 \le\frac{(2^{-r}+2^{-s})^2}{16}.
\]
Thus these matrices form a regular name with spare margin. Their squared residuals tend to zero with the explicit displayed modulus. The bases used at different stages can be unrelated.

\subsection{Represented positivity and slightly indefinite samples}
\label{sec:psdnames}

An arbitrary rational sample of a positive real matrix need not be positive. Ignoring this point would make the diagonal square-root scan invalid. We give an arithmetic correction.

Let $A_r$ be a symmetric Frobenius-regular name with $e_r=2^{-r}$. Define its positive-semidefinite contract by
\begin{equation}
 v^{\mathsf T}A_rv\ge-e_r v^{\mathsf T}v
 \qquad\text{for all numerical }r\text{ and rational vector codes }v.
 \label{eq:psdname}
\end{equation}
For ordinary real semantics this is the condition obtained from positivity of the represented matrix and a sample error at most $e_r$. It is an explicit universal arithmetic assumption, not a Boolean test on real names.

Set $B_r=A_r+e_rI$. Every rational quadratic form of $B_r$ is nonnegative by \eqref{eq:psdname}. Run the unconditional finite positivity alternative, Theorem~\ref{thm:ldl}, on $B_r$. Its negative-vector branch contradicts \eqref{eq:psdname} instantiated at the produced vector. Hence it returns a positive certificate. This is a uniform construction from the name under its proved positivity contract; it does not assume access to hidden positive rational approximants.

The correction costs at most $ne_r$ in Frobenius norm, using the convenient rational bound $\norm I\le n$ for $n\ge1$. Therefore
\[
 \bound(B_r-B_s,(n+1)(e_r+e_s)).
\]
For output stage $k$, choose a primitive recursive precision $r(k)$ satisfying
\[
 (n+1)e_{r(k)}\le\eta_k/4,
 \qquad \eta_k=2^{-2k-4}/(n+1),
\]
and apply the finite root producer to $B_{r(k)}$ with tolerance $\eta_k/4$. Then
\[
 \bound(S_k^2-S_\ell^2,(\eta_k+\eta_\ell)/2).
\]
The same mixed-basis proof establishes a regular root name. The correction vanishes with its explicit modulus, so the represented square is the original $A$. All finite approximants $S_k$ are positive, and hence the resulting name satisfies its own positivity contract.

A uniform finite magnitude bound follows either from the root certificate or from a sufficiently coarse sample and regularity. For example,
$\tr(S_k^2)=\F(S_k)$ is controlled by the diagonal entries of the nearby bounded input. Such bounds permit multiplication of the names with explicit precision demands.

\subsection{Uniqueness among all positive roots}

Theorem~\ref{thm:sqrtstable} initially assumes exact rational nonnegative diagonal certificates. Calling the output the \emph{principal} root requires uniqueness against positive roots not already supplied in that form. The following finite repair closes that gap.

\begin{lemma}[A finite uniqueness modulus]
\label{lem:rootunique}
Let $S,T$ be rational positive semidefinite matrices, $n\ge1$, with
$\bound(S,M)$ and $\bound(T,M)$ for rational $M\ge1$. Given $\eps>0$, if
\[
 \bound(S^2-T^2,\delta),\qquad
 \delta\le\frac{\eps^2}{8n},
\]
then $\bound(S-T,\eps)$.
\end{lemma}
\begin{proof}
Choose a positive rational
\[
 \alpha\le\min\left(1,\frac\eps4,
                    \frac{\eps^2}{16n(2M+1)}\right).
\]
Apply the finite diagonalizer to $S$ and $T$ at tolerance $\alpha$ and discard each off-diagonal residual. The resulting matrices $S',T'$ have exact nonnegative diagonal certificates: their diagonal entries are nonnegative quadratic forms of $S,T$. They lie within $\alpha$ of $S,T$.

For arbitrary matrices the noncommutative identity
\[
 X^2-Y^2=X(X-Y)+(X-Y)Y
\]
and finite submultiplicativity give
\[
 \norm{S'^2-S^2}\le\alpha(2M+1),\qquad
 \norm{T'^2-T^2}\le\alpha(2M+1).
\]
Here $\norm{S'},\norm{T'}\le M+1$ follows from $\alpha\le1$. Hence
\[
 \norm{S'^2-T'^2}\le\delta+2\alpha(2M+1)
 \le\frac{\eps^2}{4n}.
\]
Theorem~\ref{thm:sqrtstable} implies
$\F(S'-T')\le\eps^2/4$, that is, $\norm{S'-T'}\le\eps/2$.
Finally
\[
 \norm{S-T}\le2\alpha+\eps/2\le\eps.
\]
All approximations in this proof use a single explicitly specified rational precision. There is no appeal to an exact spectral decomposition or qualitative convergence.
\end{proof}

For two represented positive roots of the same matrix, first replace their rational samples by the positive shifted samples just constructed. Their shifts have explicit vanishing errors. Regularity supplies a common magnitude bound $M$; the displayed product identity supplies a primitive recursive modulus for their squares. At any requested output tolerance $\eps$, choose sample precisions making the square difference at most $\eps^2/(8n)$ and the sample/shift errors small enough to leave the desired margin. Apply Lemma~\ref{lem:rootunique} and the triangle inequality. This proves equality of the root names in the finite-slack sense.

We have therefore constructed the unique positive represented matrix $S$ satisfying $S^2=A$, with an explicit primitive recursive presentation. The proof covers singular positive inputs: no lower eigenvalue bound was used. The derivative of scalar square root near zero is irrelevant because the comparison has a quadratic precision demand rather than a uniform linear one.

\section{Spectral projectors at a certified gap}
\label{sec:projectors}

Ordered eigenvalue names survive collisions, but a projector selecting one side of a threshold requires that threshold to be separated from the spectrum. We give a finite condition that both states and checks the required separation.

\subsection{Finite projectors and their comparison}

Given a rational diagonalization certificate $(Q,d)$, a rational threshold $\tau$, and a rational margin $\gamma>0$, suppose
\[
 |d_i-\tau|\ge\gamma\qquad(i<n).
\]
Use the decidable rational comparison to put $p_i=1$ if $d_i<\tau$ and $p_i=0$ otherwise. Define
\[
 P=Q\diag(p)Q^{\mathsf T}.
\]
It follows by exact multiplication that $P^{\mathsf T}=P$, $P^2=P$ and $\tr(P)=\sum_i p_i$. Thus $P$ is already an exact rational orthogonal projector, even though it only approximately reduces the original input matrix.

\begin{theorem}[Projector comparison]
\label{thm:projector}
Suppose certificates $(Q,d)$ for $A$ and $(R,e)$ for $B$ have residuals $\eta,\theta$. Let both lists avoid the same $\tau$ by at least $\gamma>0$, and let $P,P'$ be their threshold projectors. If $\bound(A-B,\delta)$, then
\begin{equation}
 4\gamma^2\F(P-P')\le(\delta+\eta+\theta)^2.
 \label{eq:projector}
\end{equation}
Equivalently $\norm{P-P'}\le(\delta+\eta+\theta)/(2\gamma)$.
\end{theorem}
\begin{proof}
Write $D=\diag(d)$, $E=\diag(e)$, and $V=Q^{\mathsf T}R$. For the exact reconstructed matrices $C=QDQ^{\mathsf T}$ and $C'=RER^{\mathsf T}$,
\[
 (Q^{\mathsf T}(C-C')R)_{ij}=(d_i-e_j)V_{ij},
\]
whereas
\[
 (Q^{\mathsf T}(P-P')R)_{ij}=(p_i-p'_j)V_{ij}.
\]
If the two bits agree the second entry is zero. If they differ, one corresponding diagonal value is at most $\tau-\gamma$ and the other at least $\tau+\gamma$, so $|d_i-e_j|\ge2\gamma$. In either case,
\[
 4\gamma^2(p_i-p'_j)^2V_{ij}^2
 \le(d_i-e_j)^2V_{ij}^2.
\]
Sum and use orthogonal invariance. This gives
$4\gamma^2\F(P-P')\le\F(C-C')$.
Finally the three finite errors give
$\bound(C-C',\delta+\eta+\theta)$.
\end{proof}

This proof does not identify individual eigenvectors. It compares whole selected subspaces by a matrix observable. Eigenvalues may collide arbitrarily within the lower or upper block.

\subsection{A positive certificate for the input gap}

\begin{definition}[Rational gap certificate]
For rational symmetric $A$, a rational threshold $\tau$ and $g>0$, a gap certificate is a positive certificate for
\begin{equation}
 (A-\tau I)^2-g^2I.
 \label{eq:gapcert}
\end{equation}
It is an exact rational identity of the form $L\diag(a)L^{\mathsf T}$ with $a_i\ge0$.
\end{definition}

\begin{lemma}[A gap passes to approximate diagonals]
\label{lem:gappass}
If \eqref{eq:gapcert} is positive and $(Q,d)$ has residual at most $\eta$, then
\[
 |d_i-\tau|\ge g-\eta\qquad(i<n).
\]
In particular $\eta\le g/2$ supplies the margin $\gamma=g/2$.
\end{lemma}
\begin{proof}
For the unit column $q_i$, the gap certificate gives
\[
 \|(A-\tau I)q_i\|^2
 =q_i^{\mathsf T}(A-\tau I)^2q_i\ge g^2.
\]
The residual column in $Aq_i=d_iq_i+r_i$ has squared length at most $\eta^2$. Triangle comparison yields
\[
 \|(A-\tau I)q_i\|\le|d_i-\tau|+\eta.
\]
Squaring and comparing nonnegative rationals proves $g\le|d_i-\tau|+\eta$. All statements involve finite vector sums; the norm expressions abbreviate their squared versions.
\end{proof}

For fixed rational input with such a certificate, choose diagonalization errors
$\eta_r=g2^{-r-4}$. These are at most $g/2$. Theorem~\ref{thm:projector} with $A=B$ gives
\[
 \norm{P_r-P_s}\le2^{-r-4}+2^{-s-4}.
\]
Thus the projectors form a coherent primitive recursive matrix name. Since their exact reconstructed diagonal matrices commute with them, the residual gives
\[
 \norm{AP_r-P_rA}\le2\eta_r.
\]
To see the constant directly, multiplication by an orthogonal projector is a contraction in Frobenius norm: in its rational orthogonal coordinates it simply deletes rows or columns. Apply this to the two residual products. Finite multiplication moduli now show that the represented $P$ is symmetric, idempotent and commutes with $A$.

For represented inputs the margin conversion is also explicit. Suppose the represented gap is $G>0$, and the matrix and its rational samples have Frobenius bound $M$. Put $M_\tau=M+n|\tau|$. If a sample has input error at most $e_r$, the product identity gives squared-matrix error at most $2M_\tau e_r$. Choose a precision satisfying
\[
 2M_\tau e_r\le G^2/2
\]
and use the smaller rational gap $g=G/2$. The margin reserve is $G^2-g^2=3G^2/4$, which dominates that error. Thus the rational matrix $(A_r-\tau I)^2-g^2I$ is positive and Theorem~\ref{thm:ldl} supplies its finite certificate. Internally the premise is the corresponding universal finite quadratic-form gap contract; instantiate it at a putative negative witness from the elimination algorithm to rule that branch out. The cross-input projector comparison then handles different samples. A bare assertion that no eigenvalue equals $\tau$ is not being used as a numerical gap producer.

There is a second useful, wholly finite sufficient guard. Suppose one baseline diagonalization certificate has residual $\eta_0\le g/4$ and every diagonal value has distance at least $g$ from $\tau$. Any later certificate for the same matrix, with residual $\eta\le g/4$, has its ordered diagonal within $g/2$ of the baseline by Theorem~\ref{thm:weyl}. Its margin is at least $g/2$, and the threshold bit at each ordered index agrees with the baseline. Thus one decidable finite check can authorize all later precision demands; subsequent production does not require a new unproved existence assertion.

\subsection{Stable block counts and limits of the result}

Let $k=\tr P$ and $\ell=\tr P'$ be the integer bit counts. With $V=Q^{\mathsf T}R$,
\[
 \F(P-P')=k+\ell-2\sum_{i,j}p_ip'_jV_{ij}^2.
\]
The final sum is at most both $k$ and $\ell$, using the row and column squared sums of $V$. Hence
\[
 \F(P-P')\ge|k-\ell|.
\]
Once the comparison error is strictly less than one, the two integer counts agree. A certified gap thus supports a finite decision of the selected block size and, for a cut at zero, inertia information.

There is no unconditional rank procedure here. In the scalar family $A(t)=[t]$, rank and the zero-threshold projector change at $t=0$. Stable names for $t$ alone need not determine whether it vanishes by a finite test. Likewise the exact multiplicity partition of the ordered spectrum is not decided by the eigenvalue-name theorem. The explicit gap identifies precisely the extra data used for discrete block information.

\section{Consequences and reusable patterns}
\label{sec:applications}

\subsection{Polynomial matrix functions}

Let $p(X)=\sum_{k=0}^d a_kX^k$ have rational coefficients. If a diagonalization certificate reconstructs $C=Q\diag(d)Q^{\mathsf T}$ with residual $\eta$, and $\bound(A,M),\bound(C,M)$ for $M\ge1$, then
\begin{equation}
 \norm{p(A)-Q\diag(p(d_i))Q^{\mathsf T}}
 \le\eta\sum_{k=1}^d k|a_k|M^{k-1}.
 \label{eq:polycalc}
\end{equation}
Indeed $p(C)=Qp(\diag(d))Q^{\mathsf T}$ by exact finite multiplication. The noncommuting telescoping identity
\[
 A^k-C^k=\sum_{j=0}^{k-1} A^j(A-C)C^{k-1-j}
\]
follows by cancellation of successive terms. Repeated left or right multiplication by $A,C$ costs at most a factor $M$ per multiplication. Thus each summand has Frobenius bound $M^{k-1}\eta$, including the endpoint cases in which one of the powers is the identity. Summing proves \eqref{eq:polycalc}. There is no erroneous Frobenius bound $\norm I=1$ in this argument: identity multiplication simply leaves the matrix unchanged.

The coefficient constant in \eqref{eq:polycalc} is a primitive recursive precision transformer. It gives stable matrix-polynomial observations despite changing diagonalization witnesses. Trace changes by at most $n$ times the Frobenius error, since each diagonal entry has at most that error. Therefore finite sums of polynomial eigenvalue values also recover the corresponding trace identities with explicit budgets.

An analytic matrix function can be handled if supplied with rational polynomial approximants, a primitive recursive uniform approximation modulus on a known interval, and the relevant finite error proof. The framework propagates those inputs. It does not manufacture the scalar approximation theorem or its modulus.

\subsection{Singular values and absolute values}

For a rational rectangular matrix $V$, its Gram matrix $A=V^{\mathsf T}V$ is positive because
\[
 x^{\mathsf T}Ax=\sum_j(Vx)_j^2\ge0.
\]
Finite elimination therefore supplies the positive certificate used above. The ordered eigenvalue names of $A$ and bounded scalar square-root searches give $n$ right-Gram singular-value names for an $m\times n$ input, with multiplicities and zeros, without any choice of singular vectors. This convention retains the additional forced zeros when $n>m$, whereas a common convention lists only $\min(m,n)$ values. The elementary scalar inequality
\[
 |\sqrt u-\sqrt v|^2\le|u-v|\qquad(u,v\ge0)
\]
is interpreted through the same rational square-root certificates; its proof is $(s-t)^2\le|s^2-t^2|$ plus the finite-slack argument.

For two finite matrices $V,W$,
\[
 V^{\mathsf T}V-W^{\mathsf T}W
 =V^{\mathsf T}(V-W)+(V-W)^{\mathsf T}W.
\]
Hence a common magnitude bound and Frobenius input error give an explicit Gram-matrix error, then Theorem~\ref{thm:weyl} gives eigenvalue errors, and the square-root comparison gives singular-value errors. This is a concrete composition of three certificate interfaces.

For symmetric represented $A$, define its matrix absolute value by
$|A|=(A^2)^{1/2}$. The preceding finite constructions supply its name without making a sign decision on any eigenvalue. The expressions $(|A|+A)/2$ and $(|A|-A)/2$ are likewise computable matrix names. By comparison with the scalar functions in arbitrarily fine diagonalization certificates, both are positive and their difference is $A$; the finite polynomial and root moduli justify this interpretation. For an explicit positivity budget, let $C=QDQ^{\mathsf T}$ approximate $A$ within $\eta$, with both Frobenius bounds at most $M$, and put $T=Q\diag(|d_i|)Q^{\mathsf T}$. Then $T^2=C^2$ and $\norm{T^2-A^2}\le2M\eta$. If $S$ is a positive spectral root certificate for $A^2$ of residual $\rho$, Theorem~\ref{thm:sqrtstable} gives
\[
 \F(S-T)\le n(\rho+2M\eta).
\]
Whenever $\sigma^2\ge n(\rho+2M\eta)$, each $(S\pm A)/2$ lies within $(\sigma+\eta)/2$ of the exactly positive matrix $Q\diag((|d_i|\pm d_i)/2)Q^{\mathsf T}$. This supplies the claimed positivity contract at every finite tolerance. A support projector or a polar factor at zero would require a separate gap contract and is not inferred from these names.

\subsection{What has been resolved}

The existing frameworks already distinguish approximate witnesses from coherent points. This report supplies a concrete method for resolving some cases where insisting on coherent witnesses would be too strong:
\begin{itemize}
\item \emph{Colliding roots in a structured subclass.} Real symmetric characteristic polynomials have ordered, multiplicity-preserving root names with a cross-input comparison theorem. This goes beyond their root supports.
\item \emph{Unstable bases.} The intrinsic square root and gap projectors have coherent names even if every finite diagonalizer selects a different basis. A global eigenbasis is not an intermediate premise.
\item \emph{Hidden finite-dimensional existence.} Rational elimination gives explicit witnesses for subspace intersection and positive certificates. Their uniform dimension bounds are part of the proof.
\item \emph{Loss of positivity during approximation.} The $A_r+e_rI$ correction and finite positivity alternative give an explicit repair with an error budget, rather than presuming that rational samples preserve a closed condition.
\item \emph{Unexplained convergence.} The spectral computation has an additive rational energy invariant and explicit fuel. The resulting name coherence follows from a separate certificate comparison theorem.
\end{itemize}

The general lesson is operational. First ask for a finite output at one tolerance. Next determine the observation that should be independent of this finite choice. Prove a comparison estimate for arbitrary certificates and distinct nearby inputs. Only then define the infinite mathematical object by its regular observations. If a discrete selection is required, isolate the numerical separation needed to make that selection stable.

\subsection{A small exact example}

For
\[
 A=\begin{pmatrix}1&1\\1&1\end{pmatrix}
\]
the first pivot has $h=0,b=1$. Taking $(c,s)=(3/5,4/5)$ gives
\[
 Q=\frac15\begin{pmatrix}3&-4\\4&3\end{pmatrix},\qquad
 Q^{\mathsf T}AQ=
 \begin{pmatrix}49/25&-7/25\\-7/25&1/25\end{pmatrix}.
\]
The original off-diagonal energy is $2$; the new energy is $98/625$. After sorting, the diagonal observations are $(1/25,49/25)$. Both orthogonality identities are exact. Further finite stages improve the residual to any prescribed rational tolerance, while the ordered observations name $(0,2)$.

The positive square root is represented without requiring the irrational eigenvector normalization to appear as an exact rational witness. Its finite certificates satisfy the explicitly checked equation $S_r^2\approx A$. At a threshold strictly between $0$ and $2$, the lower spectral projector is also a coherent matrix observation. At threshold zero there is no positive separation from the spectrum, so the same gap theorem makes no claim.

% Research fragment for the main report. All repository references use the
% pinned revision c39c388d4f87220ab9e758ae457520bbc0d74341.
\section{The boundary between certificates and global witnesses}
\label{sec:certificate-boundary}

The audited arithmetic packets \cite{faltings,siegel,mordell,roth,prime} make an essential distinction that a general
embedding framework must preserve. A finite computation can check a
candidate point, a rational identity, or a finite formal derivation.
It cannot automatically certify that all points of an infinite search
space have already been found. The latter is a mathematical assertion
about every possible point. There is also a distinction between a proof
that some bound exists and a certificate which contains a particular
numerical bound and certifies its correctness.

\subsection{An exact impossibility theorem for complete bound certificates}

Fix an effective enumeration of deterministic programs. Let $H(e,t)$ mean
that program $e$, run on input $e$, halts for the first time at positive
step $t$. Bounded simulation and a bounded check of the preceding steps
make $H$ primitive recursive. Determinism implies
\[
 H(e,s)\wedge H(e,t)\ \longrightarrow\ s=t.
\]
Thus every fiber has the uniform cardinality bound one. Set
\[
 \operatorname{Bound}(e,B)
   \quad\Longleftrightarrow\quad
 \forall t\bigl(H(e,t)\longrightarrow t\le B\bigr).
\]

\begin{theorem}[No complete effective system of numerical bound certificates]
\label{thm:no-complete-bound-certificates}
There is no total computable predicate $C(e,B,w)$ satisfying both of the
following properties in the standard natural numbers:
\begin{enumerate}
\item \emph{Soundness:} $C(e,B,w)$ implies
      $\operatorname{Bound}(e,B)$.
\item \emph{Completeness:} for every $e$ there are $B,w$ with $C(e,B,w)$.
\end{enumerate}
The same impossibility holds for a computably enumerable acceptance
relation. In particular, it holds for primitive recursive checkers.
\end{theorem}

\begin{proof}
Suppose first that $C$ is decidable. Given $e$, enumerate all pairs
$(B,w)$ in a fixed effective order and test $C(e,B,w)$. Completeness
guarantees that this search terminates. Its running time need not have
a primitive recursive bound; ordinary recursive termination is already
enough for the argument. Let $B_e$ be the bound in the first accepted
certificate.

Simulate program $e$ on input $e$ for $B_e$ steps. If it halts, answer
yes. If it does not halt within those steps, answer no. Soundness makes
the second answer correct: any later first halting time $t$ would satisfy
$t>B_e$, contrary to $\operatorname{Bound}(e,B_e)$. Hence this procedure
would be a total decision algorithm $D$ for diagonal halting.

For completeness, the diagonal contradiction is immediate. Construct a
program $P$ which, on input $n$, runs $D(n)$ and loops forever when the
answer is yes, but halts when the answer is no. Let $p$ be a code of $P$.
If $D(p)$ answers yes, then $P(p)$ loops; if it answers no, then $P(p)$
halts. Either answer contradicts the asserted correctness of $D$.

If acceptance is only computably enumerable, dovetail the acceptance
computations for all pairs $(B,w)$. Completeness again supplies a pair
whose acceptance will eventually be observed, and the remainder of the
argument is unchanged.
\end{proof}

The theorem concerns certificates which explicitly supply a numerical
bound. It does not say that a classical theory cannot prove
$\forall e\exists B\operatorname{Bound}(e,B)$: classical case analysis
on whether $e(e)$ halts proves that assertion. A finite classical proof
of the existential sentence does not by itself provide a numerical
witness. Nor does the theorem say that any particular family of
algebraic curves has this obstruction. The halting fibers are a
counterexample to a general certificate-completeness principle, not
an encoding of curves.

This yields a practical requirement for the framework. Every use of
certificates should identify separately the finite acceptance relation,
its proved soundness, and the existence or construction of accepted
certificates on the intended input class. A sound checker may be a useful
artifact even when completeness is unavailable. It cannot serve as a
substitute for the missing completeness theorem.

\subsection{The precise intuitionistic witnessing theorem}

The strong witnessing assertion used in the Siegel packet is supported
by a primary proof-theoretic reference. Salehi defines $\mathsf{iPRA}$
using all primitive recursive defining equations and atomic induction,
and states the following as Fact~3.5 on page~57:
\[
 \begin{gathered}
 T\vdash\forall x\exists y\,A(x,y)\\
 \Longrightarrow\quad
 \text{some primitive recursive symbol }f\text{ satisfies}\\
 T\vdash\forall x\,A(x,f(x)).
 \end{gathered}
\]
for $T=\mathsf{iPRA}$ or intuitionistic $\Sigma_1$ induction in the
primitive recursive language. The matrix $A$ is not restricted to a
decidable or quantifier-free formula. See \cite{salehi}.
This is an external metatheorem, not an internal reflection axiom and
not a result formalized by the present files. Atomic and quantifier-free
induction are interderivable here by the Boolean representations of
quantifier-free formulas.

Ferreira and Marques give a related formulation for intuitionistic
$I\Sigma_1$, with correctness in its primitive-recursive definitional
extension; their author manuscript states the theorem on printed
page~14:
\cite{ferreira}.
The Faltings packet cites that more conservative formulation and does
not claim an extracted correctness proof in $\mathsf{iPRA}$ itself.
Salehi's exact statement supports the stronger conclusion when its
specified $\mathsf{iPRA}$ hypothesis is met.

Here is the premise-management step in full. Suppose $G(x)$ is decidable
in $\mathsf{iPRA}$ and
\[
 \mathsf{iPRA}\vdash
 \forall x\bigl(G(x)\longrightarrow\exists y\,A(x,y)\bigr).
\]
Fix $x$. If $G(x)$, the displayed theorem supplies $y$ and $A(x,y)$,
and therefore $G(x)\to A(x,y)$. If $\neg G(x)$, choose $y=0$;
the implication $G(x)\to A(x,0)$ follows by contradiction. Decidability
of $G$ proves these cases exhaustive. Generalization gives
\[
 \mathsf{iPRA}\vdash
 \forall x\exists y\bigl(G(x)\longrightarrow A(x,y)\bigr).
\]
The witnessing theorem now supplies one primitive recursive $f$ and
an $\mathsf{iPRA}$ proof of
\[
 \forall x\bigl(G(x)\longrightarrow A(x,f(x))\bigr).
\]
This does not use unrestricted independence of premise. Its essential
hypothesis is the proved decidability of $G$.

Apply this observation with
$A(c,B)=\forall p(P(c,p)\to h(p)\le B)$. A uniform positive
$\mathsf{iPRA}$ theorem on decidably admissible curve codes gives a
primitive recursive height bound, even though $A$ contains an unbounded
universal quantifier. If $c$ includes a certificate, the extracted
function is allowed to depend on that certificate. A further producer
is necessary before the input can be reduced to raw curve coefficients.
The implication gives a necessary computational consequence; absence
of a known bound is not a nonprovability theorem.

\subsection{What the audited arithmetic packets still require}

The following statements are relative to the pinned repository revision,
not a claim that all mathematics in the packets has been verified.

\paragraph{Faltings and Siegel.}
The positive list and height versions require a global height bound.
The packets prove bounded enumeration after such a bound is supplied
and distinguish a cardinality bound from a height bound. The Siegel
packet already proves the finite covering last step of a quantitative
Subspace Theorem route: at most $M$ subspaces, each meeting the curve
in at most $D$ relevant points, give a cardinality bound $MD$. Its
missing ingredients include a finite-tuple approximation theorem,
the requisite curve-function presentations, lifting and valuation
estimates, and their arithmetic proofs. The present general framework
does not supply those ingredients.

\paragraph{Mordell--Weil.}
The supplied-cover result is substantially more explicit than an
abstract invocation of descent. For an integral short Weierstrass model
$y^2=x^3+Ax+B$, it gives a primitive recursive list of all halves of
any supplied rational point by a rational-root scan of a quartic. From
a supplied complete list of representatives modulo doubling it derives
an explicit height contraction and generator list. Pointwise
decidability of divisibility by two does not construct that complete
quotient-covering list. General production of that list remains an
unfilled obligation.

\paragraph{Roth.}
The quadratic case and the Liouville range are treated by explicit
polynomial separation. General cutoff effectivity is left open. The
packet already includes the cardinality-to-empty-window construction
using $N+1$ successive intervals, so that observation should not be
presented as a new contribution.

\paragraph{Chebotarev.}
The packet explicitly separates effective analytic estimates from their
translation to $\mathsf{iPRA}$. Its conditional density-modulus lemma
is supplied with integer thresholds. The still missing inputs include
number-field and prime-ideal algorithms with arithmetic correctness,
Dedekind/Hecke explicit formulas, quantitative zero-free estimates,
and certified primitive recursive constants. A possible exceptional
zero alone is not asserted to prevent effective Chebotarev.

\paragraph{Siegel--Walfisz and Bombieri--Vinogradov.}
The report isolates uniform constants and thresholds before the
unbounded size variable. Its elementary reduction from effective
Bombieri--Vinogradov witnesses to logarithmic-power Siegel--Walfisz
estimates shows why averaging alone does not erase the issue. It also
retains the aggregate rational approximation error $QX2^{-s}$ when
summing over $Q$ moduli. Neither finite evaluation of prime counts nor
a finite Gram-matrix certificate supplies the missing uniform analytic
estimates.

\subsection{Lean targets for the certificate boundary}

The computability obstruction is a metatheoretic theorem, naturally
formalized in Lean's ordinary metatheory rather than as a new
$\mathsf{iPRA}$ axiom. A clean dependency chain is:
\begin{enumerate}
\item Fix a concrete machine model and a primitive recursive bounded
simulation, with the first-halt predicate and its uniqueness proof.
\item Define $\operatorname{Bound}$ and a certificate structure with
an acceptance predicate and semantic soundness theorem.
\item Given totality of accepted pairs, define the partial-recursive
least accepted pair; prove its totality from the hypothesis. No claim
that this minimization is primitive recursive is needed.
\item Compose it with bounded simulation and prove that it decides
diagonal halting, using soundness exactly when the simulation reports
no halt by the supplied bound.
\item Apply the chosen machine model's diagonal undecidability
theorem, or formalize the short self-application contradiction above.
\end{enumerate}
For iPRA integration, only the finite simulator, first-halt uniqueness,
and finite acceptance checks are object-language kernels. The
impossibility theorem and the external witnessing theorem belong in
the metatheory and must not be added as arithmetic axioms.


\section{Lean definitions, checked glue, and the completion path}
\label{sec:lean}

\subsection{What is compiled}

The archive contains a small reproducible Lake project pinned to Lean 4.32.0, Mathlib revision
\code{81a5d257c8e410db227a6665ed08f64fea08e997}, and the three unchanged ConciseIPRA modules \code{Syntax.lean}, \code{Deduction.lean}, and \code{PR.lean} from the source revision. The source hashes and upstream license are included.

\code{ObjectObservation.lean} compiles five generic inference combinators whose conclusions have the actual embedded \code{Prf PRAAx} type. They are not assertions about the standard natural numbers substituted for derivations. Their roles are:
\begin{enumerate}
\item \code{boundedInvariantAt}: use permitted quantifier-free induction and instantiate at a supplied numerical fuel term;
\item \code{loopProducesCertificate}: combine the terminal invariant with a supplied object-level extraction implication;
\item \code{observableOfIndependentCertificates}: apply a supplied stability implication to two derived certificates;
\item \code{observableFromLocalStability}: discharge two local certificate assumptions and apply the resulting implication;
\item \code{twoLoopsStableObservable}: combine two finite-loop derivations and a local observation comparison.
\end{enumerate}

Their axiom audit reports \code{propext} and \code{Quot.sound}, with no \code{sorryAx} or \code{Classical.choice}. These are Lean metatheoretic dependencies and do not add axioms to the embedded object calculus.

\code{SpectralBlueprint.lean} separately gives compiler-checked host definitions of finite rational matrices, sums, algorithms, certificates and target propositions. Definitions ending in \code{Target} name unproved propositions; they are neither axioms nor proved theorems. The structure \code{UnprovedMatrixObligations} lists them and has no supplied instance. Its rational implementation inherits \code{Classical.choice} transitively through Std, in addition to \code{propext} and \code{Quot.sound}. This host dependency is explicitly recorded. The file is not used as a derivation inside $\ipra$; the numerical reification below must have its own arithmetic proofs.

\subsection{The exact upstream interfaces}

The primitive recursive syntax is \code{PR : Nat -> Type}, generated by zero, successor, projection, composition and primitive recursion. Its term type is \code{PTm n}, generated by variables and application of a PR symbol. The step arguments for primitive recursion are ordered \emph{time, previous value, parameters}. Formula syntax and intuitionistic natural deduction are supplied by \code{Syntax.lean} and \code{Deduction.lean} \cite{prsource,deductionsource}.

For $\varphi$ over $n+1$ free variables, the existing rule \code{PRA.qfInd} takes
\[
 \varphi.\mathrm{QF},\qquad
 \Gamma\vdash_{\mathrm{PRAAx}}\varphi[0],\qquad
 \varphi::\mathrm{shift}(\Gamma)\vdash_{\mathrm{PRAAx}}\varphi[Sv_0]
\]
and returns $\Gamma\vdash_{\mathrm{PRAAx}}\forall v_0\,\varphi$.
The new \code{boundedInvariantAt} simply follows this with universal elimination at the explicit fuel term. The context may contain quantified input contracts; the induction formula itself must remain quantifier-free.

The final existing bridge is
\[
 (\Gamma\vdash_{\mathrm{PRAAx}}\varphi)
 \ \longleftrightarrow\
 \operatorname{NPrf}(\Gamma.\mathrm{map}\,\operatorname{NFm.ofPRA},
                            \operatorname{NFm.ofPRA}(\varphi)),
\]
named \code{pra_iff_native} \cite{nativesource}. Applying its forward direction is the last step after the desired embedded proof has been built. It is not a compiler from arbitrary Mathlib theorems to $\ipra$.

\subsection{The concrete host definitions}

The blueprint uses
\begin{lstlisting}
abbrev Vec (n : Nat) := Fin n -> Rat
abbrev Mat (n : Nat) := Fin n -> Fin n -> Rat

def sumQ : (n : Nat) -> (Fin n -> Rat) -> Rat
  | 0, _ => 0
  | n + 1, f => f 0 + sumQ n (fun i => f i.succ)
\end{lstlisting}
From this recursive sum it defines transpose, multiplication, diagonal matrices, reconstruction $QDQ^{\mathsf T}$, conjugation, squared Frobenius norm, off-diagonal energy, matrix-vector multiplication and quadratic forms. There is no real completion in these definitions.

Certificate data and validity are deliberately separated. \code{DiagonalData n} has a rational basis and vector of values. \code{DiagonalCertificate A eta w} states symmetry, nonnegative tolerance, both orthogonality identities, sortedness, exact agreement with the transformed diagonal, and the squared residual bound. \code{PositiveData} has a general rational factor and nonnegative weights, not an orthogonal factor; otherwise generic positive rational matrices would be excluded. \code{RootCertificate} asserts a positive root candidate and a squared residual. \code{ProjectorCertificate} asserts symmetry, idempotence and a commutator residual.

The file also defines the exact two-branch \code{pivotRotation}, plane rotations, finite first-maximum pivot selection, the frozen \code{jacobiStep}, primitive-recursive \code{jacobiRun}, the explicit \code{jacobiFuel}, and diagonal sorting. A harmless total fallback in the sorted-list lookup is eliminated by the usual sorting-length proof. The producer's mathematical correctness is the unproved proposition \code{DiagonalizationFuelTarget}, not inferred merely from type checking its definition.

\subsection{Numerical encodings for the object layer}

One explicit encoding is sufficient to make the path unambiguous. Use signed-integer codes $0\mapsto0$, positive $z\mapsto2z$, and negative $-z\mapsto2z-1$ for $z>0$. A convenient positive pairing is
\[
 \operatorname{pair}(a,b)=2^a(2b+1).
\]
Its projections use bounded search for the exponent of two and integer division. The code zero has fixed default projections. For a positive pair the second component is smaller than the pair code, and all projections are primitive recursive. Use pair-coded signed numerator and denominator-minus-one for rationals, with zero again given the default value $0/1$.

Encode a list by its length and iterated positive pairs, and a matrix by its dimension and a row-major list of length $n^2$. A validity checker verifies exactly the requested number of cells. Every raw operation is totalized on invalid codes. Cross-multiplication proves rational equality and order; normalization can be omitted from the logical encoding. For the computational bit bound use the reduced or common-denominator representation specified in Section~\ref{sec:grid}.

The fixed functions to reify as \code{PR} syntax include rational arithmetic; finite sums and products; array access/update; pivot scanning; plane-rotation update; state iteration; rational ceiling; final permutation; the positivity stack and its reconstruction; greatest-square bounded search; and threshold masks. The ordinary PR closure theorem applies to each fixed definition. A \code{PR} code for each function and a correctness relation to its host meaning are construction tasks, not supplied automatically by the Python program or Lean's termination checker.

\subsection{The proof obligations in dependency order}

\begin{longtable}{@{}P{0.25\textwidth}P{0.43\textwidth}P{0.26\textwidth}@{}}
\toprule
Target & Main proof and induction data & Required output \\
\midrule\endhead
Arithmetic representation & Signed order; positive denominators; cross-products; finite Boolean folds. Induct on numerical recursion counters with all parameters fixed. & Object equations and reflection lemmas for each fixed checker. \\
Finite matrix algebra & Expand bounded sums; prove reindexing, transpose and multiplication laws; finite Cauchy--Schwarz by the sum-of-squares identity. & Encoded identities and squared inequalities. \\
Pivot improvement and energy identity & Rational branch split; guarded denominator clearing; the two scalar inequalities and preservation of outside-pivot squared sums. & Local object proof for one finite update. \\
Finite diagonalization & The Booleanized invariant \eqref{eq:invariant}, exact orthogonality and conjugation, maximum-pivot scan, final sorting. & \code{Prf PRAAx} of producer correctness at the explicit fuel. \\
Rational intersection & Row-operation stack, pivot/free columns, back substitution, and nonzero free coordinate; at most $n$ stages. & A PR rational intersection witness with its finite checker proof. \\
Sorted spectral comparison & Use that witness, sorted coordinate bounds and finite Cauchy--Schwarz; clear the positive $x^{\mathsf T}x$ denominator. & Object cross-input comparison for arbitrary certificates. \\
Positivity targets & Bounded Schur elimination and reverse reconstruction, exact congruence invariant, negative-vector alternative. & Producer correctness for weighted positive factors or negative witnesses. \\
Root producer and mixed-basis comparison & Bounded square-grid search, exact orthogonal reconstruction, weighted Cauchy--Schwarz. & Root residual and comparison; then the finite uniqueness modulus. \\
Projector targets & Bits square to themselves; mixed-basis separation; positive gap matrix; finite row/column squared sums. & Projector identity, residual, gap margin, and cross-input comparison. \\
Name schemas & Instantiate the finite comparison at computed precisions, under fixed universal input-name contracts. & Arithmetic finite-observation theorems for every output precision. \\
Native transfer & Apply the already proved \code{pra_iff_native.mp}. & Native-core derivation of the translated formula. \\
\bottomrule
\end{longtable}

Every variable-size statement in the matrix rows is expressed by a checker on a finite code. For example, the orthogonality component is the Boolean product over $i,j<n$ of the equality tests for the corresponding finite dot products. The invariant does not quantify over all rational vectors or all future precisions. The universal quadratic-form theorem is established by proving the finite identity with its vector as a free numerical parameter and generalizing afterward. Similarly, a matrix name's universal regularity or positivity statement remains a fixed context assumption instantiated at needed numerical arguments.

For determinant multiplication, expand the two permutation sums. Terms indexed by a repeated intermediate column cancel in pairs under the transposition of the repeated positions; the remaining terms are indexed by permutations and factor into the two determinant sums. This is a bounded involution and reindexing argument. The principal-minor coefficient identity follows by choosing which diagonal $X$ terms occur in the permutation expansion. Its finite binomial and factorial counts have their own prefix-count invariants. These supply the elementary algebra required by Lemma~\ref{lem:coeff} without importing a semantic spectral theorem.

\subsection{What would count as a completed formalization}

A completed application must contain actual numerical PR definitions, object-level correctness proofs for the finite arithmetic kernels, the producer and comparison derivations, and the final instantiated theorem. Supplying an object formula with an axiom naming the desired theorem would not meet this criterion. Neither would proving only its semantic interpretation in Lean, even if that ambient proof is constructive.

The supplied compiled combinators already enforce the right interface: their missing local estimates are arguments of type \code{Prf PRAAx}, not arbitrary Lean propositions. They cannot silently absorb a host spectral theorem. The blueprint and reference programs make the remaining mathematical obligations concrete and testable, but the main matrix-specific embedded proofs remain to be constructed.

\section{Artifacts, verification, and remaining work}
\label{sec:validation}

\subsection{Executable exact arithmetic}

The Python reference uses integers and \code{fractions.Fraction} throughout. It contains no floating-point tolerance tests in certificate acceptance. \code{spectral.py} implements both the exact-ratio and finite-grid diagonalizers, together with independent certificate checking. \code{observables.py} implements the positive-factor/negative-vector alternative, root certificates and threshold projectors. The programs are references for the finite constructions, not proof checkers for the whole theory $\ipra$.

The included verification scripts check the following finite cases:
\begin{center}
\begin{tabular}{lr}
\toprule
Check family & Cases \\
\midrule
Scalar pivot inequalities, both rational branches & 600 \\
Full three-dimensional energy identities & 3 \\
Complete exact-ratio diagonalization certificates & 9 \\
Complete finite-grid diagonalization certificates & 4 \\
Positive rational factorization certificates & 9 \\
Negative quadratic-form witnesses & 5 \\
Positive square-root residual certificates & 8 \\
Mixed-basis square-root comparison instances & 3 \\
Sorted spectral comparisons & 3 matrix pairs / 8 coordinates \\
Exact rational projector certificates & 5 \\
Cross-precision projector comparison & 1 \\
\bottomrule
\end{tabular}
\end{center}
Additional negative tests reject invalid input contracts and a deliberately corrupted root certificate. The examples cover dimensions zero and one, indefinite matrices, repeated spectra, a split collision, singular positive matrices, nonsquare scalar roots and dense matrices of dimensions three and four.

In one dense $4\times4$ grid example, tolerance $1/16$ was achieved in 21 rotations; the largest denominator of an entry of $Q$ had 658 bits. These are reproducible facts about the included example, not asymptotic performance claims. The uniform proof is the energy argument and its explicit fuel, independent of these runs.

\subsection{Lean and provenance}

The standalone Lean build checks the embedded inference combinators and the host blueprint definitions separately. The axiom logs distinguish the two layers, including the transitive host-rational \code{Classical.choice} dependency. No main spectral theorem is reported as proved by the kernel.

The provenance manifest identifies the pinned repository revision, exact inspected source paths, and hashes of the locally retrieved source text. The report's bibliography supplies human-readable primary references. The source audit is targeted; the claims of other proof packets are not promoted to verified theorems merely because they were read.

\subsection{Theorem status ledger}

\begin{longtable}{@{}P{0.36\textwidth}P{0.27\textwidth}P{0.31\textwidth}@{}}
\toprule
Result & Mathematical status here & Machine status here \\
\midrule\endhead
Stable-observation composition & Full elementary inference argument. & Compiled generic embedded proof combinators. \\
Rational approximate diagonalization and grid variant & Full finite proof with explicit fuel and Booleanizable invariants. & Exact reference computations; host algorithm and target types checked; object proof pending. \\
Sorted multiplicity-preserving spectral names & Full cross-certificate comparison and coefficient-factorization argument. & Host target types and finite comparison tests; object proof pending. \\
Rational positivity alternative & Full bounded elimination argument. & Exact producer/checker and negative witnesses; formal proof pending. \\
Principal positive square root & Full residual, stability, sample-correction and uniqueness arguments. & Exact certificates and comparison tests; formal proof pending. \\
Gap-dependent spectral projector & Full finite comparison and gap-production arguments. & Exact projectors and comparison checks; formal proof pending. \\
No complete effective numerical-bound certificates for arbitrary singleton fibers & Full metatheoretic diagonal proof. & Formalization plan only. \\
General Siegel/Faltings, general Mordell--Weil quotient covers, general effective Roth, full prime-distribution targets & Not proved here; remaining input theorems identified. & No new claimed formalization. \\
\bottomrule
\end{longtable}

The most direct next formalization target is the scalar pivot lemma, followed by the energy identity and fuel invariant. These close the producer side without any real analysis. The rational intersection and spectral comparison then close the stable-observation side. Together they would yield the first substantive embedded matrix theorem of this package. The square-root and projector developments reuse those same finite kernels.

\begin{thebibliography}{99}

\bibitem{repoindex}

ProveIt, Primitive recursive arithmetic: directory overview.

\url{https://github.com/VladimirReshetnikov/ProveIt/blob/c39c388d4f87220ab9e758ae457520bbc0d74341/Logic/PrimitiveRecursiveArithmetic/README.md}.

\bibitem{packetindex}

ProveIt, Delivered iPRA proof packets: intake and status index.

\url{https://github.com/VladimirReshetnikov/ProveIt/blob/c39c388d4f87220ab9e758ae457520bbc0d74341/Logic/PrimitiveRecursiveArithmetic/docs/proofs/README.md}.

\bibitem{contracts}

Finite contracts for mathematics in iPRA, representations section.

\url{https://github.com/VladimirReshetnikov/ProveIt/blob/c39c388d4f87220ab9e758ae457520bbc0d74341/Logic/PrimitiveRecursiveArithmetic/docs/proofs/foundations/ipra_framework/paper/sections/03_representations.tex}.

\bibitem{observations}

Finite observations and uniform constructions, analysis section.

\url{https://github.com/VladimirReshetnikov/ProveIt/blob/c39c388d4f87220ab9e758ae457520bbc0d74341/Logic/PrimitiveRecursiveArithmetic/docs/proofs/foundations/ipra_finite_observation/tex/analysis.tex}.

\bibitem{aggregates}

Finite demands, root supports and norm certificates, aggregate methods.

\url{https://github.com/VladimirReshetnikov/ProveIt/blob/c39c388d4f87220ab9e758ae457520bbc0d74341/Logic/PrimitiveRecursiveArithmetic/docs/proofs/foundations/ipra_finite_demand/sections/03_aggregates.tex}.

\bibitem{rootsupport}

Finite demands, root supports and norm certificates, root-set methods.

\url{https://github.com/VladimirReshetnikov/ProveIt/blob/c39c388d4f87220ab9e758ae457520bbc0d74341/Logic/PrimitiveRecursiveArithmetic/docs/proofs/foundations/ipra_finite_demand/sections/04_root_sets.tex}.

\bibitem{factorization}

Finite observations and uniform constructions, uniform factorization.

\url{https://github.com/VladimirReshetnikov/ProveIt/blob/c39c388d4f87220ab9e758ae457520bbc0d74341/Logic/PrimitiveRecursiveArithmetic/docs/proofs/foundations/ipra_finite_observation/tex/factorization.tex}.

\bibitem{coherent}

Coherent finite methods, finite refinement presentations.

\url{https://github.com/VladimirReshetnikov/ProveIt/blob/c39c388d4f87220ab9e758ae457520bbc0d74341/Logic/PrimitiveRecursiveArithmetic/docs/proofs/foundations/ipra_coherent_methods_2026-11-15/paper/sections/03_refinement.tex}.

\bibitem{localglobal}

Local-to-global constructions, object iteration.

\url{https://github.com/VladimirReshetnikov/ProveIt/blob/c39c388d4f87220ab9e758ae457520bbc0d74341/Logic/PrimitiveRecursiveArithmetic/docs/proofs/foundations/ipra_local_global/paper/03_object_iteration.tex}.

\bibitem{rotations}

Constructive Euler--Rodrigues packet, overview and logical boundaries.

\url{https://github.com/VladimirReshetnikov/ProveIt/blob/c39c388d4f87220ab9e758ae457520bbc0d74341/Logic/PrimitiveRecursiveArithmetic/docs/proofs/geometry/euler_rodrigues_ipra/README.md}.

\bibitem{faltings}

Faltings in iPRA: finiteness formulations and proof-theoretic analysis.

\url{https://github.com/VladimirReshetnikov/ProveIt/blob/c39c388d4f87220ab9e758ae457520bbc0d74341/Logic/PrimitiveRecursiveArithmetic/docs/proofs/diophantine-number-theory/faltings_ipra/faltings_ipra.tex}.

\bibitem{siegel}

Siegel in iPRA: logic and global-boundary analysis.

\url{https://github.com/VladimirReshetnikov/ProveIt/blob/c39c388d4f87220ab9e758ae457520bbc0d74341/Logic/PrimitiveRecursiveArithmetic/docs/proofs/diophantine-number-theory/siegel_ipra/paper/02_logic.tex}.

\bibitem{mordell}

Constructive Mordell--Weil, supplied-cover elliptic descent.

\url{https://github.com/VladimirReshetnikov/ProveIt/blob/c39c388d4f87220ab9e758ae457520bbc0d74341/Logic/PrimitiveRecursiveArithmetic/docs/proofs/diophantine-number-theory/mordell_weil_ipra/paper/03_elliptic.tex}.

\bibitem{roth}

Diophantine approximation in iPRA: Roth cutoffs and finite windows.

\url{https://github.com/VladimirReshetnikov/ProveIt/blob/c39c388d4f87220ab9e758ae457520bbc0d74341/Logic/PrimitiveRecursiveArithmetic/docs/proofs/diophantine-number-theory/diophantine_ipra_2026-11-15/paper/02_roth.tex}.

\bibitem{prime}

Prime distribution in iPRA: Siegel--Walfisz and Bombieri--Vinogradov targets.

\url{https://github.com/VladimirReshetnikov/ProveIt/blob/c39c388d4f87220ab9e758ae457520bbc0d74341/Logic/PrimitiveRecursiveArithmetic/docs/proofs/analytic-number-theory/prime_distribution_ipra/paper/sections/07_sw_bv.tex}.

\bibitem{prsource}

ConciseIPRA: primitive recursive symbols, terms, axioms and induction.

\url{https://github.com/VladimirReshetnikov/ProveIt/blob/c39c388d4f87220ab9e758ae457520bbc0d74341/Logic/PrimitiveRecursiveArithmetic/ConciseIPRA/Lean/ConciseIPRA/PR.lean}.

\bibitem{deductionsource}

ConciseIPRA: intuitionistic natural deduction.

\url{https://github.com/VladimirReshetnikov/ProveIt/blob/c39c388d4f87220ab9e758ae457520bbc0d74341/Logic/PrimitiveRecursiveArithmetic/ConciseIPRA/Lean/ConciseIPRA/Deduction.lean}.

\bibitem{nativesource}

ConciseIPRA: native bi-interpretation and transfer.

\url{https://github.com/VladimirReshetnikov/ProveIt/blob/c39c388d4f87220ab9e758ae457520bbc0d74341/Logic/PrimitiveRecursiveArithmetic/ConciseIPRA/Lean/ConciseIPRA/NativeBiInterpretation.lean}.

\bibitem{salehi}

Saeed Salehi, \emph{Decidable Formulas of Intuitionistic Primitive Recursive Arithmetic}, Reports on Mathematical Logic 36 (2002), 55--61, especially Fact 3.5, p.~57.

\url{https://rml.tcs.uj.edu.pl/rml-36/36-SALEHI.pdf}.

\bibitem{ferreira}

Fernando Ferreira and Ant\'onio Marques, \emph{Extracting Algorithms from Intuitionistic Proofs}, author manuscript, especially the intuitionistic $I\Sigma_1$ discussion on printed p.~14.

\url{https://webpages.ciencias.ulisboa.pt/~fjferreira/extracting.pdf}.

\bibitem{jacobinotes}

University of Pennsylvania course text, \emph{Numerical Analysis}, excerpt, sections 2.7--2.9 on eigenvalues and the Jacobi method. Used for historical/methodological context; the rational progress proof here is given independently.

\url{https://www2.math.upenn.edu/~deturck/m320/text/part8.pdf}.

\end{thebibliography}

\end{document}
